% Tapis de Quillen, with Quillen's offprint and preprint (cote 162-5, typed part): archive cote 162-5, site page 3 (PDF page 4).
% typescript. Draft compared: none.
% AI-transcribed by Claude Opus 5.5 (claude-opus-5-5, adaptive thinking, effort high), Message Batches API, batch t162_full (msgbatch_01UAh3zyWPd454AeDkhnmQ33), 2026-11-13; not yet checked by eye.
\scanpage{162-5/3}
\mspage{\hand{2}}

On a aussi un foncteur évident
\[ T: (\mathrm{Ssimpl}) \longrightarrow (\mathrm{Cat}) , \]
en associant à tout ensemble semi-simplicial X la catégorie $T(X) = \hand{\Delta}_{/X}$
des simplexes sur X, dont l'ensemble des objets est la réunion d\unsure{i|a}sjointe
des $X_n$ .... (c'est une catégorie fibrée sur la catégorie $\hand{\Delta}$ des simplexes
types, à fibres les catégories discrètes définies par les $X_n$). Ceci posé
Quillen prouve que pour tout X, ST(X) est isomorphe canoniquement à X da
% NOTE: the typed line runs off the right edge of the sheet after "da"
la catégorie homotopique construite avec (Ssimpl), et que pour toute C,
la catégorie TS(C) est \inserted[typed]{canoniquement} ``homotopiquement équivalente à C'' \xout{\unsure{dansxlaxxxxx}} i.e
\inserted[typed]{\unsure{canoniquement}} isomorphe à C dans la catégorie quotient de (Cat) obtenue en inversant le
% NOTE: the typed line runs off the right edge of the sheet after "le"
foncteurs qui sont des homotopismes. \xout{Ilxenxrésulte} Ces isomorphismes son
% NOTE: the typed line runs off the right edge of the sheet after "son"
fonctoriels en X. Il en résulte formellement qu'un morphisme $f:X\rightarrow Y$ dans
(Ssimpl) est un homotopisme si et seulement si il en est ainsi de $T(f):T($
% NOTE: the typed line runs off the right edge of the sheet after "T("
$\rightarrow T(Y)$, d'où des foncteurs $S'$: $(\mathrm{Cat})' \longrightarrow (\mathrm{Ssimpl})'$ et $T'$:$(\mathrm{Ssimpl})' \longrightarrow (\mathrm{C}$
% NOTE: the typed line runs off the right edge of the sheet after "(C"
entre les catégories ``homotopiques'', \inserted[typed]{construites avec (Cat) resp. (Ssimpl),} qui sont quasi-inverses l'un de l'au
re.

De plus, Quillen construit un isomorphisme canonique et fonctoriel
dans (Cat)' entre C et la catégorie opposée $C^{o}$ , ou ce qui revient au
même, un isomorphisme canonique et fonctoriel dans (Ssimpl)' entre S(C) e
% NOTE: the typed line runs off the right edge of the sheet after "e"
$S(C^{o})$. \xout{\unsure{Celaxluixpermetxalors,xpourx}} La définition est telle que le
foncteur induit sur les systèmes locaux sur C transforme le foncteur
contravariant F sur C, transformant toute flêche en flêche inversible, en
le foncteur covariant (i.e. contravariant sur $C^{o}$) ayant mêmes valeurs sur
les objets, et obtenu sur les flêches en remplaçant $F(u)$ par $F(u)^{-1}$; en
d'autres termes, l'effet de l'homotopisme de Quillen sur les groupoides
fondamentaux est l'isomorphisme évident entre les groupoides fondamentaux
de C et de $C^{o}$ ,\xout{\unsure{quixestxlixdxxx}} compte tenu que le deuxième est l'opposé d
% NOTE: the typed line runs off the right edge of the sheet after "d"

\margin{\hand{3}}
